% section: チェビシェフ多項式：三角関数を多項式に翻訳する

\section{チェビシェフ多項式：三角関数を多項式に翻訳する}

\subsection{同じ漸化式が二つの姿を持つ}

\[
  a_{n+1}=2x a_n-a_{n-1}
\]
を考える.
$x$ を固定すれば二階線形漸化式であり,特性方程式は
\[
  r^2-2xr+1=0
\]
である.

$-1\leq x\leq1$ のとき,$x=\cos\theta$ と書ける.
すると,
\[
  \cos((n+1)\theta)
  =
  2\cos\theta\cos(n\theta)-\cos((n-1)\theta)
\]
だから,$\cos(n\theta)$ は同じ漸化式を満たす.

ここで,「三角関数の値を計算する問題」と「漸化式を解く問題」が同じ式を持っている.
この一致を,変数 $x$ の多項式として保存したくなる.

\subsection{チェビシェフ多項式の定義と証明}

第一種チェビシェフ多項式 $T_n(x)$ を,
\[
  T_0(x)=1,\qquad T_1(x)=x
\]
および
\[
  T_{n+1}(x)=2xT_n(x)-T_{n-1}(x)
\]
で定める.

初めのいくつかは,
\[
  \begin{aligned}
    T_0(x) & =1,          \\
    T_1(x) & =x,          \\
    T_2(x) & =2x^2-1,     \\
    T_3(x) & =4x^3-3x,    \\
    T_4(x) & =8x^4-8x^2+1
  \end{aligned}
\]
である.

漸化式の右辺が多項式だけでできているので,すべての $T_n(x)$ は多項式になる.
一方,
\[
  T_n(\cos\theta)=\cos(n\theta)
\]
が成り立つ.

実際,$n=0,1$ では定義から明らかである.
$n$ と $n-1$ について成り立つと仮定すれば,
\[
  \begin{aligned}
    T_{n+1}(\cos\theta)
     & =2\cos\theta T_n(\cos\theta)-T_{n-1}(\cos\theta) \\
     & =2\cos\theta\cos(n\theta)-\cos((n-1)\theta)      \\
     & =\cos((n+1)\theta)
  \end{aligned}
\]
となる.

三角関数を使って書かれた量を,$x$ の多項式に翻訳できた.
例えば,
\[
  \cos(3\theta)=4\cos^3\theta-3\cos\theta
\]
は,
\[
  T_3(x)=4x^3-3x
\]
という多項式の公式になる.

\subsection{三種類の領域：振動,境界,増大}

特性方程式
\[
  r^2-2xr+1=0
\]
の根は,
\[
  r=x\pm\sqrt{x^2-1}
\]
である.

\begin{itemize}
  \item $|x|<1$ なら根は複素数で,単位円上にある.$T_n(x)$ は振動する.
  \item $x=1$ または $x=-1$ では根が重なり,境界的な振る舞いになる.
  \item $|x|>1$ なら実根の絶対値が異なり,一方の根が支配的になる.
\end{itemize}

$x=\cos\theta$ の代わりに $x=\cosh t$ と置くと,
\[
  T_n(\cosh t)=\cosh(nt)
\]
となる.
三角関数の振動と双曲線関数の増大が,同じ多項式の異なる領域で現れている.

これは,特性根の絶対値が数列の増加を支配するという,線形漸化式の基本原理を別の言葉で見たものである.

\subsection{第二種チェビシェフ多項式}

第一種だけでなく,\[
  U_0(x)=1,\qquad U_1(x)=2x
\]
および
\[
  U_{n+1}(x)=2xU_n(x)-U_{n-1}(x)
\]
で定まる第二種チェビシェフ多項式も考えられる.
こちらは,
\[
  U_n(\cos\theta)
  =\frac{\sin((n+1)\theta)}{\sin\theta}
\]
を満たす.

第一種が余弦を記録するのに対して,第二種は正弦を記録する.
同じ漸化式でも初期値を変えると,異なる多項式族が現れる.
初期値は単なる付属条件ではなく,解の空間の中からどの解を選ぶかを決めている.

\subsection{直交性と近似への入口}

チェビシェフ多項式が特別なのは,三角関数を表すからだけではない.
余弦の直交性から,
\[
  \int_0^\pi\cos(m\theta)\cos(n\theta)\,d\theta=0
  \qquad(m\neq n)
\]
が得られる.

$x=\cos\theta$ と変数変換すれば,
\[
  \int_{-1}^{1}
  \frac{T_m(x)T_n(x)}{\sqrt{1-x^2}}\,dx=0
  \qquad(m\neq n)
\]
となる.
この性質を直交性という.

直交する多項式を使うと,複雑な関数を
\[
  c_0T_0(x)+c_1T_1(x)+c_2T_2(x)+\cdots
\]
のように展開して近似できる.
ここからフーリエ級数や数値計算,最良近似の理論へ進むことができる.

本書で必要なのは近似理論の完成ではない.
一つの漸化式が,三角関数を経由して,多項式の直交性や関数近似へつながるという道筋を見つけることである.

\clearpage
